% TODO increase workload sizes so we can compare really large workloads
% TODO build time in OpenCL

\subsection{Results for NVIDIA GeForce GTX 970}
The benchmark program using the Adam optimizer and the \ac{GEMM} workload was first executed
on an NVIDIA GeForce GTX 970 with 8 GB of VRAM. For this benchmark, the following results
were obtained for data transfer time and computation time using CUDA and OpenCL.

\begin{table}[th]
	\caption{Data transfer time and computation time across all batch indices}
	\label{tab:gtx970results}
	\centering
	\begin{NiceTabular}{lrrr}
		\CodeBefore
		\rowcolors{2}{gray!10}{white}
		\Body
		\textbf{Metric} & \textbf{CUDA} & \textbf{OpenCL} & \textbf{Perf. Ratio} \\
		Average data-transfer time (ms) & 4.06 & 4.43 & ? \% \\
		Average computation time (ms) & 2.45 & 0.77 & ? \% \\
	\end{NiceTabular}
\end{table}

Table \ref{tab:gtx970results} presents the benchmark results obtained with the specified
hyperparameters. The data transfer time was averaged across all measurements performed for
both OpenCL and CUDA. The results show that CUDA is 9\% faster in terms of data transfer time
on the NVIDIA GeForce GTX 970. Similarly, the computation time was averaged over all runs,
resulting in CUDA being 66\% faster than OpenCL.

\begin{figure}[th]
	\begin{center}
		\includegraphics[width=16cm, height=8cm]{nvidia_geforce_gtx_970_compute_and_transfer_batch_line_normalized}
	\end{center}
	\caption{Results of running the GEMM workload using the Adam optimizer}
	\label{fig:pltgtx970}
	\centering
\end{figure}

Figure \ref{fig:pltgtx970} shows the average computation time per \verb|batch_index| and
also includes the average loss per \verb|batch_index| for both frameworks. This illustrates
the efficiency of both frameworks while also providing an indication of their reliability.
The figure shows that the loss decreases over time, which matches the expected behavior
during optimization. Furthermore, the figure highlights the results summarized in
Table \ref{tab:gtx970results}. This data shows that on data transfer time CUDA provides 
better results while for computation OpenCL provides better results on the GTX 970 representing
here the Maxwell architecture from NVIDIA.

\subsection{Results for RTX 5000}
\subsection{Results for A-100}